\input{../tex-inputs/latex-leseansicht-vorspann}

\section[Theodor von Sosnosky an Arthur Schnitzler, 24. 1. 1901]{L04228 Theodor von Sosnosky an Arthur Schnitzler, 24. 1. 1901}
\nopagebreak\mylabel{L04228v}
\rehead{ }\normalsize\beginnumbering\briefempfaengerindex{Schnitzler, Arthur@\textsc{Schnitzler, Arthur}!zzzSosnosky, Theodor von@\emph{von Theodor von Sosnosky}!1901-01-241@{24. 1. 1901}|(be}
\toendnotes[C]{\smallbreak\pagebreak[2]}
\correspDesc{Versand  durch Theodor von Sosnosky am 24. 1. 1901 in Kremsmünster
\newline{}Erhalt  durch Arthur Schnitzler im Zeitraum [25. 1. 1901
                  – 29. 1. 1901?] in Wien}\toendnotes[C]{\smallbreak}
\Standort{DLA, A:Schnitzler, HS.1985.1.4640.}
\physDesc{Brief, 1 Blatt, 3 Seiten, 1908 Zeichen
\newline{}Handschrift: blaue Tinte, deutsche Kurrent
\newline{}Schnitzler: mit rotem Buntstift eine Unterstreichung }\toendnotes[C]{\smallbreak}
\pstart{}{\pb}Sehr geehrter Herr Doktor!\pend\vspace{0.5em}
\pstart
           Dieser Tage wird Ihnen von \textsc{Cotta}\orgindex{J.G. Cotta’sche Buchhandlung Nachfolger@J.G. Cotta’sche Buchhandlung Nachfolger|pw} meine Anthologie\pwindex{deutsche Lyrik des 19. Jahrhunderts. Eine poetische Revue@\emph{Die deutsche Lyrik des 19. Jahrhunderts. Eine poetische Revue}|pwv} zugehen, in der
               auch Sie vertreten ſind. Ich hatte vor, Ihnen perſönlich ein\strikeout{es}{ }\introOben{}Exemplar\introOben{}\pwindex{deutsche Lyrik des 19. Jahrhunderts. Eine poetische Revue@\emph{Die deutsche Lyrik des 19. Jahrhunderts. Eine poetische Revue}|pwv} zu überſenden, um mich ſowohl für Ihren freundlichen \label{K_L04228-1v}\edtext{Beitrag\pwindex{Schnitzler, Arthur 15.\,5.\,1862 Wien – 21.\,10.\,1931 ebd.@\textsc{Schnitzler, Arthur} (15.\,5.\,1862 Wien – 21.\,10.\,1931 ebd.), \emph{Schriftsteller, Mediziner}!Ohnmacht@\strich\emph{Ohnmacht}|pwv}\pwindex{Schnitzler, Arthur 15.\,5.\,1862 Wien – 21.\,10.\,1931 ebd.@\textsc{Schnitzler, Arthur} (15.\,5.\,1862 Wien – 21.\,10.\,1931 ebd.), \emph{Schriftsteller, Mediziner}!Anfang vom Ende@\strich\emph{Anfang vom Ende}|pwv}\pwindex{Schnitzler, Arthur 15.\,5.\,1862 Wien – 21.\,10.\,1931 ebd.@\textsc{Schnitzler, Arthur} (15.\,5.\,1862 Wien – 21.\,10.\,1931 ebd.), \emph{Schriftsteller, Mediziner}!Wie wir so still@\strich\emph{Wie wir so still}|pwv}}{\lemma{\textnormal{\emph{Beitrag}}}\Cendnote{\textnormal{Zu Sosnoskys\pwindex{Sosnosky, Theodor von 4.\,1.\,1866 Budapest – 3.\,2.\,1943 Wien@\textsc{Sosnosky, Theodor von} (4.\,1.\,1866 Budapest – 3.\,2.\,1943 Wien), \emph{Schriftsteller}|pwk} Antologie \emph{Die deutsche Lyrik des 19. Jahrhunderts. Eine poetische
                     Revue}\pwindex{deutsche Lyrik des 19. Jahrhunderts. Eine poetische Revue@\emph{Die deutsche Lyrik des 19. Jahrhunderts. Eine poetische Revue}|pwk} hatte Schnitzler die drei
                  Gedichte \emph{Ohnmacht}\pwindex{Schnitzler, Arthur 15.\,5.\,1862 Wien – 21.\,10.\,1931 ebd.@\textsc{Schnitzler, Arthur} (15.\,5.\,1862 Wien – 21.\,10.\,1931 ebd.), \emph{Schriftsteller, Mediziner}!Ohnmacht@\strich\emph{Ohnmacht}|pwk}, \emph{Anfang vom Ende}\pwindex{Schnitzler, Arthur 15.\,5.\,1862 Wien – 21.\,10.\,1931 ebd.@\textsc{Schnitzler, Arthur} (15.\,5.\,1862 Wien – 21.\,10.\,1931 ebd.), \emph{Schriftsteller, Mediziner}!Anfang vom Ende@\strich\emph{Anfang vom Ende}|pwk} und \emph{Wie
                     wir so still}\pwindex{Schnitzler, Arthur 15.\,5.\,1862 Wien – 21.\,10.\,1931 ebd.@\textsc{Schnitzler, Arthur} (15.\,5.\,1862 Wien – 21.\,10.\,1931 ebd.), \emph{Schriftsteller, Mediziner}!Wie wir so still@\strich\emph{Wie wir so still}|pwk} beigesteuert.}}}\label{K_L04228-1} als für die Spende Ihres charmanten Buches\pwindex{Schnitzler, Arthur 15.\,5.\,1862 Wien – 21.\,10.\,1931 ebd.@\textsc{Schnitzler, Arthur} (15.\,5.\,1862 Wien – 21.\,10.\,1931 ebd.), \emph{Schriftsteller, Mediziner}!Reigen. Zehn Dialoge@\strich\emph{Reigen. Zehn Dialoge}|pwv} zu revanchiren; aber der
                  Verlag\orgindex{J.G. Cotta’sche Buchhandlung Nachfolger@J.G. Cotta’sche Buchhandlung Nachfolger|pwv} teilte mir mit, daß
               er ohnehin an \uline{alle} im Buche\pwindex{deutsche Lyrik des 19. Jahrhunderts. Eine poetische Revue@\emph{Die deutsche Lyrik des 19. Jahrhunderts. Eine poetische Revue}|pwv} vertretenen Autoren, ſofern ſie noch
               am Leben sind, Frei-Exemplare\pwindex{deutsche Lyrik des 19. Jahrhunderts. Eine poetische Revue@\emph{Die deutsche Lyrik des 19. Jahrhunderts. Eine poetische Revue}|pwv}
               verſende. Da dies wahrſcheinlich{ }ſchon geſchehen iſt, ſo ſende ich ſelber kein Exemplar\pwindex{deutsche Lyrik des 19. Jahrhunderts. Eine poetische Revue@\emph{Die deutsche Lyrik des 19. Jahrhunderts. Eine poetische Revue}|pwv} an Sie ab, ſondern
               begnüge mich mit dieſem Briefe.\pend
           
\pstart
           Bei dieſer Gelegenheit will ich auch bemerken, daß meine Abſicht, Ihr Buch\pwindex{Schnitzler, Arthur 15.\,5.\,1862 Wien – 21.\,10.\,1931 ebd.@\textsc{Schnitzler, Arthur} (15.\,5.\,1862 Wien – 21.\,10.\,1931 ebd.), \emph{Schriftsteller, Mediziner}!Reigen. Zehn Dialoge@\strich\emph{Reigen. Zehn Dialoge}|pwv} zu beſprechen, keineswegs leeres
                Geflunker geweſen {\pb}iſt, wie es den Anſchein hat, da noch i{\geminationm}er keine Rezenſion erſchienen iſt; ich habe vielmehr in
               einem längeren Eſſai »Jung Wien\pwindex{Sosnosky, Theodor von 4.\,1.\,1866 Budapest – 3.\,2.\,1943 Wien@\textsc{Sosnosky, Theodor von} (4.\,1.\,1866 Budapest – 3.\,2.\,1943 Wien), \emph{Schriftsteller}!Jung-Wien«. Studienblätter@\strich\emph{»Jung-Wien«. Studienblätter}|pw}« ihre \uline{geſa{\geminationm}te} litterariſche
               Thätigkeit kritiſch zu ſchildern verſucht und dabei ſehr \label{K_L04228-2v}\edtext{auſführlich von »\textsc{Reigen\pwindex{Schnitzler, Arthur 15.\,5.\,1862 Wien – 21.\,10.\,1931 ebd.@\textsc{Schnitzler, Arthur} (15.\,5.\,1862 Wien – 21.\,10.\,1931 ebd.), \emph{Schriftsteller, Mediziner}!Reigen. Zehn Dialoge@\strich\emph{Reigen. Zehn Dialoge}|pw}}«}{\lemma{\textnormal{\emph{ausführlich von »Reigen«}}}\Cendnote{\textnormal{Der Essay wurde in drei Teilen gegeben: der erste Teil war am 20. 1. 1901 erschienen und widmete sich je zur Hälfte Hermann Bahr\pwindex{Bahr, Hermann 19.\,7.\,1863 Linz – 15.\,1.\,1934 München@\textsc{Bahr, Hermann} (19.\,7.\,1863 Linz – 15.\,1.\,1934 München), \emph{Schriftsteller, Kritiker}|pwk} und Schnitzler Werk, der zweite Teil erschien am 26. 1. 1901 und enthielt eine lange Passage über den \emph{Reigen}\pwindex{Schnitzler, Arthur 15.\,5.\,1862 Wien – 21.\,10.\,1931 ebd.@\textsc{Schnitzler, Arthur} (15.\,5.\,1862 Wien – 21.\,10.\,1931 ebd.), \emph{Schriftsteller, Mediziner}!Reigen. Zehn Dialoge@\strich\emph{Reigen. Zehn Dialoge}|pwk} neben der Besprechung von Werken von Hugo von Hofmannsthal\pwindex{Hofmannsthal, Hugo von 1.\,2.\,1874 Wien – 15.\,7.\,1929 Rodaun@\textsc{Hofmannsthal, Hugo von} (1.\,2.\,1874 Wien – 15.\,7.\,1929 Rodaun), \emph{Schriftsteller}|pwk} und Richard Beer-Hofmann\pwindex{Beer-Hofmann, Richard 11.\,7.\,1866 Wien – 26.\,9.\,1945 New York City@\textsc{Beer-Hofmann, Richard} (11.\,7.\,1866 Wien – 26.\,9.\,1945 New York City), \emph{Schriftsteller}|pwk}. Der dritte Teil wurde am 3. 2. 1901 publiziert und behandelte Peter Altenberg\pwindex{Altenberg, Peter 9.\,3.\,1859 Wien – 8.\,1.\,1919 ebd.@\textsc{Altenberg, Peter} (9.\,3.\,1859 Wien – 8.\,1.\,1919 ebd.), \emph{Schriftsteller}|pwk}, Felix Dörmann\pwindex{Dörmann, Felix 29.\,5.\,1870 Wien – 26.\,10.\,1928 ebd.@\textsc{Dörmann, Felix} (29.\,5.\,1870 Wien – 26.\,10.\,1928 ebd.), \emph{Schriftsteller}|pwk}, Felix Salten\pwindex{Salten, Felix 6.\,9.\,1869 Budapest – 8.\,10.\,1945 Zürich@\textsc{Salten, Felix} (6.\,9.\,1869 Budapest – 8.\,10.\,1945 Zürich), \emph{Schriftsteller, Journalist, Chefredakteur}|pwk} und Ludwig Wolff\pwindex{Wolff, Ludwig 7.\,3.\,1876 Bielsko-Biała – nach 1958 Vereinigte Staaten von Amerika [USA]@\textsc{Wolff, Ludwig} (7.\,3.\,1876 Bielsko-Biała – nach 1958 Vereinigte Staaten von Amerika [USA]), \emph{Schriftsteller, Dramaturg}|pwk} (\emph{»Jung-Wien«. Studienblätter von Theodor v. Sosnosky}\pwindex{Sosnosky, Theodor von 4.\,1.\,1866 Budapest – 3.\,2.\,1943 Wien@\textsc{Sosnosky, Theodor von} (4.\,1.\,1866 Budapest – 3.\,2.\,1943 Wien), \emph{Schriftsteller}!Jung-Wien«. Studienblätter@\strich\emph{»Jung-Wien«. Studienblätter}|pwk}. In: Beilage zur
                        \emph{Norddeutschen Allgemeinen Zeitung}\pwindex{Norddeutsche Allgemeine Zeitung@\emph{Norddeutsche Allgemeine Zeitung}|pwk},
                     Jg. 40, Nr. 17a, 20. 1. 1901,
                     S. [1]–2, Nr. 20, 26. 1. 1901, S. [1] und Nr. 29a, 3. 2. 1901, S. [1]–2).}}}\label{K_L04228-2} geſprochen. Dieſer Aufſatz\pwindex{Sosnosky, Theodor von 4.\,1.\,1866 Budapest – 3.\,2.\,1943 Wien@\textsc{Sosnosky, Theodor von} (4.\,1.\,1866 Budapest – 3.\,2.\,1943 Wien), \emph{Schriftsteller}!Jung-Wien«. Studienblätter@\strich\emph{»Jung-Wien«. Studienblätter}|pwv} ist im Juli an die »\uline{Nord-deutſche Allgem.\pwindex{Norddeutsche Allgemeine Zeitung@\emph{Norddeutsche Allgemeine Zeitung}|pw}}« abgegangen, die ſich geneigt erklärt hat, ihn zu bringen.\pend
           
\pstart
           Bis heute hab’ ich nichts gehört! Doch werde ich den Aufſatz\pwindex{Sosnosky, Theodor von 4.\,1.\,1866 Budapest – 3.\,2.\,1943 Wien@\textsc{Sosnosky, Theodor von} (4.\,1.\,1866 Budapest – 3.\,2.\,1943 Wien), \emph{Schriftsteller}!Jung-Wien«. Studienblätter@\strich\emph{»Jung-Wien«. Studienblätter}|pwv} ſchließlich \uline{doch} irgendwo herausbringen, falls er dort etwa nicht
               erſcheinen ſollte.\pend
           
\pstart
           Eine Beſprechung des Buches\pwindex{Schnitzler, Arthur 15.\,5.\,1862 Wien – 21.\,10.\,1931 ebd.@\textsc{Schnitzler, Arthur} (15.\,5.\,1862 Wien – 21.\,10.\,1931 ebd.), \emph{Schriftsteller, Mediziner}!Reigen. Zehn Dialoge@\strich\emph{Reigen. Zehn Dialoge}|pwv} –
               ohne Nennung Ihres Namens, – alſo um{ }ſo pikanter für die Leſer – hatte ich einer
               \label{K_L04228-3v}\edtext{litterarischen Revue\pwindex{Sosnosky, Theodor von 4.\,1.\,1866 Budapest – 3.\,2.\,1943 Wien@\textsc{Sosnosky, Theodor von} (4.\,1.\,1866 Budapest – 3.\,2.\,1943 Wien), \emph{Schriftsteller}!Literatur. Aus der Bilanz der Belletristik@\strich\emph{Literatur. Aus der Bilanz der Belletristik}|pwv}}{\lemma{\textnormal{\emph{litterarischen Revue}}}\Cendnote{\textnormal{Th. v. Sosnosky\pwindex{Sosnosky, Theodor von 4.\,1.\,1866 Budapest – 3.\,2.\,1943 Wien@\textsc{Sosnosky, Theodor von} (4.\,1.\,1866 Budapest – 3.\,2.\,1943 Wien), \emph{Schriftsteller}|pwk}: \emph{Literatur. Aus der Bilanz der Belletristik}\pwindex{Sosnosky, Theodor von 4.\,1.\,1866 Budapest – 3.\,2.\,1943 Wien@\textsc{Sosnosky, Theodor von} (4.\,1.\,1866 Budapest – 3.\,2.\,1943 Wien), \emph{Schriftsteller}!Literatur. Aus der Bilanz der Belletristik@\strich\emph{Literatur. Aus der Bilanz der Belletristik}|pwk}. In: \emph{Neues Wiener Abendblatt}\pwindex{Neues Wiener Abendblatt@\emph{Neues Wiener Abendblatt}|pwk}. Abend-Ausgabe des »\emph{Neuen Wiener Tagblatt}\pwindex{Neues Wiener Tagblatt@\emph{Neues Wiener Tagblatt}|pwk}«, Jg. 35, Nr. 34, 4. 2. 1901, S. 4.}}}\label{K_L04228-3}, die demnächst im »\textsc{N. Wiener Tagblatt\pwindex{Neues Wiener Tagblatt@\emph{Neues Wiener Tagblatt}|pw}}« erſcheinen wird, einverleibt, erregte jedoch – \label{K_L04228-4v}\edtext{\textsc{sub rosa}}{\lemma{\textnormal{\emph{sub rosa}}}\Cendnote{\textnormal{lateinisch: unter der Rose, unter dem Siegel der Verschwiegenheit (übertr.)}}}\label{K_L04228-4} bemerkt! – Anſtoß damit und mußte den
                diesbezüglichen {\pb}Teil ſtreichen. Das liebe Publikum fühlt ſich ja ſo leicht
               in ſeiner imaginären Keuſchheit verletzt.\pend
           
\pstart
           Wenn ich meine Zuſage alſo bisher nicht gehalten hab, ſo bitte das freundlichſt nicht
               auf mein Schuld-\textsc{Conto} zu ſetzen.\pend
           
\pstart
           Indem ich die mir gütigſt zur Verfügung geſtellten 3 Gedichte\pwindex{Schnitzler, Arthur 15.\,5.\,1862 Wien – 21.\,10.\,1931 ebd.@\textsc{Schnitzler, Arthur} (15.\,5.\,1862 Wien – 21.\,10.\,1931 ebd.), \emph{Schriftsteller, Mediziner}!Ohnmacht@\strich\emph{Ohnmacht}|pwv}\pwindex{Schnitzler, Arthur 15.\,5.\,1862 Wien – 21.\,10.\,1931 ebd.@\textsc{Schnitzler, Arthur} (15.\,5.\,1862 Wien – 21.\,10.\,1931 ebd.), \emph{Schriftsteller, Mediziner}!Anfang vom Ende@\strich\emph{Anfang vom Ende}|pwv}\pwindex{Schnitzler, Arthur 15.\,5.\,1862 Wien – 21.\,10.\,1931 ebd.@\textsc{Schnitzler, Arthur} (15.\,5.\,1862 Wien – 21.\,10.\,1931 ebd.), \emph{Schriftsteller, Mediziner}!Wie wir so still@\strich\emph{Wie wir so still}|pwv} zu
               anderweitiger Verwendung wieder mit verbindlichſtem Danke zurückſtelle, zeichne ich,
               geehrter Herr Doktor, {\\[\baselineskip]}mit vorzüglicher{\\[\baselineskip]}Hochachtung{\\[\baselineskip]}\spacefill\mbox{Theodor vSosnosky}\pend
           \leftskip=0em{}
\pstart
           \noindent{}\uline{derzeit}: \textsc{Kremsmünster},
                     \textsc{Ob-Oesterreich}\oindex{Kremsmünster@\textbf{Kremsmünster}, \emph{Verwaltungsgebiet}|pw}\pend
           
\pstart
           24./I. 1901.\pend
           \selectlanguage{ngerman}\endnumbering\briefempfaengerindex{Schnitzler, Arthur@\textsc{Schnitzler, Arthur}!zzzSosnosky, Theodor von@\emph{von Theodor von Sosnosky}!1901-01-241@{24. 1. 1901}|)be}\mylabel{L04228h}
\begin{anhang}
\end{anhang}\newcommand{\dateiname}{L04228}\newcommand{\titel}{Theodor von Sosnosky an Arthur Schnitzler, 24. 1. 1901}\newcommand{\editorInnen}{Herausgegeben von Jahnke, SelmaMüller, Martin Anton}\input{../tex-inputs/latex-leseansicht-abspann}
